%% ============================================================
%% EXTENDED PHYSICS -- split 1 of 3 from the bundled file.
%% Drop-in for <<KEEP:pressure-drop>> <<KEEP:temp-prop>> <<KEEP:delay>>.
%% Insert BEFORE the existing drafted paragraph in that subsection (the one reconciling
%% the 0.03 % pump-energy share against the Wilo capacity and the pandapipes check),
%% which this file deliberately does NOT duplicate.
%% Source: v1 lines 967-1133, restructured per phenomenon rather than per variant.
%% ============================================================

Three phenomena are added above the node-resolved baseline $\Lone$, each on its own rung so
that no comparison changes more than one thing: temperature propagation along pipes
($\Ltwo$), trunk pressure drop and the pumping power it implies ($\Lthree$), and transport
delay ($\Lsix$). A fourth object, the nonlinear reference $\NLref$, retains all three in
native form and is used for forward evaluation and bounded re-solves rather than as a rung
of the ladder. Piecewise-linear approximations use SOS2 machinery with $K = 3$ segments
throughout. Both the supply and the return circuit are resolved at every level from
$\Lthree$ upward.

\paragraph{Pressure drop and pumping power ($\Lthree$).}
Friction in pipe $p$ follows Darcy--Weisbach, relating the pressure drop $\Delta p_p$ [Pa]
to the square of the mass flow $\dot{m}_p$ [kg/s]:
%
\begin{equation}
  \Delta p_{p} = f_D\, \frac{L_p}{D_p}\,
    \frac{8\,\dot{m}_p^2}{\rho \pi^2 D_p^4},
  \label{eq:darcy}
\end{equation}
%
with $f_D$ the Darcy friction factor, $D_p$ [m] the inner diameter and $\rho$ [kg/m$^3$]
the fluid density. Because the supply and return temperatures are prescribed by the
heating curve, the mass flow is linear in the pipe heat flow $Q_{p,t}$ [MW]:
%
\begin{equation}
  \dot{m}_{p,t} = \frac{Q_{p,t}}{
    c_p \bigl(T_{\mathrm{sup}}^{\mathrm{nom}}(t)
    - T_{\mathrm{ret}}^{\mathrm{nom}}(t)\bigr)},
  \label{eq:massflow}
\end{equation}
%
so that $\Delta p_{p,t} = R_p\, Q_{p,t}^2$ with $R_p$ [Pa/MW$^2$] the hydraulic resistance.
In $\Lthree$ this quadratic is approximated piecewise-linearly over $K$ equal segments with
SOS2 weights $w_{p,t,k} \in [0,1]$,
%
\begin{equation}
  \Delta p_{p,t}^{\mathrm{PWL}} = \sum_{k=1}^{K}
    m_k^{(p)}\, w_{p,t,k},
  \label{eq:pressure_pwl}
\end{equation}
%
whose maximum approximation error is $\varepsilon^{\max}_{\Delta p} = R_p \bar{Q}_p^2/(4K^2)$,
about $2.8$\,\% of the peak pressure drop at $K = 3$ (Appendix~\ref{app:pwl}). In $\NLref$
the quadratic enters natively.

Pumping power is the hydraulic work summed over pipes and divided by the pump efficiency,
%
\begin{equation}
  P_t^{\mathrm{pump}} = \frac{1}{\eta_{\mathrm{pump}}}
    \sum_{p \in \mathcal{P}}
    \frac{\dot{m}_{p,t}\, \Delta p_{p,t}}{\rho},
  \label{eq:pump_power}
\end{equation}
%
which is cubic in $Q_{p,t}$ and is piecewise-linearised directly in $\Lthree$; in $\NLref$
it is carried by the auxiliary $W_{p,t} = Q_{p,t}^2$ entering as $Q_{p,t} W_{p,t}$.

\paragraph{Temperature propagation ($\Ltwo$).}
Where $\Lone$ computes pipe losses from fixed nominal temperatures, $\Ltwo$ tracks the
outlet temperature as the fluid travels. In steady state it decays exponentially from the
inlet temperature towards the ground temperature $T_{\mathrm{gr}}$:
%
\begin{equation}
  T_{p,t}^{\mathrm{out}} = T_{\mathrm{gr}}(t)
    + \bigl(T_{n_1,t}^{\mathrm{sup}} - T_{\mathrm{gr}}(t)\bigr)
    \cdot \underbrace{\exp\!\left(
      \frac{-U_p L_p}{\dot{m}_{p,t}\, c_p}
    \right)}_{\displaystyle\phi_p(\dot{m}_{p,t})}.
  \label{eq:temp_prop}
\end{equation}
%
The decay factor $\phi_p \in (0,1]$ is the fraction of the inlet temperature excess
surviving to the outlet: near unity at high flow, near zero as flow approaches zero. Two
nonlinearities follow: the exponential itself and the bilinear product
$T_{n_1,t}^{\mathrm{sup}} \phi_p$. The exponential is approximated by a PWL in $\dot{m}$
shared by every level that uses it, so this component of the error is common and cancels in
the comparisons. The bilinear product enters natively in $\NLref$; in $\Ltwo$ it is
linearised by a first-order Taylor expansion about the nominal operating point,
%
\begin{align}
  T_{p,t}^{\mathrm{out}} &\approx T_{\mathrm{gr}}(t)
    + \bigl(T^{\mathrm{nom}}(t) - T_{\mathrm{gr}}(t)\bigr)
      \phi_{p,t}^{\mathrm{PWL}} \notag \\
    &\quad + \phi_p^{\mathrm{nom}}
      \bigl(T_{n_1,t}^{\mathrm{sup}} - T^{\mathrm{nom}}(t)\bigr),
  \label{eq:temp_prop_lin}
\end{align}
%
with the neglected second-order cross-term bounded by
$|\Delta T_{\max}||\Delta\phi_{\max}| < 2.5$\,K (Appendix~\ref{app:taylor_derivation}).

One property of this formulation determines how it is reported. Freeing the node temperatures
makes the programme non-convex and, absent a hydraulic penalty on the resulting flow
increase, degenerate: the optimiser reduces the supply temperature uniformly to its lower
bound, collapsing the loss term rather than reproducing the physical decay. In the solve it
does so for about $91$\,\% of hours, against a mean supply temperature of $85.9$\,°C at the
neighbouring levels. The solved objective of $\Ltwo$ is therefore not a meaningful fidelity
measurement and is not used; the effect of temperature propagation is quantified
instead through the forward evaluator of Section~\ref{subsec:evaluator}, which applies the
native exponential to a fixed schedule and cannot exploit the degeneracy. $\Ltwo$ remains in
the ladder as the formulation whose linearisation error is measured, not as a source of
dispatch cost.

\paragraph{Transport delay ($\Lsix$).}
At finite velocity, heat injected at the source at time $t$ arrives downstream later. The
nominal travel time through pipe $p$, with cross-section $A_p = \pi D_p^2/4$, is
%
\begin{equation}
  \tau_p = \frac{L_p\, \rho\, A_p}{\dot{m}_p^{\mathrm{nom}}},
  \label{eq:delay_time}
\end{equation}
%
discretised to $k_p = \mathrm{round}(\tau_p/\Delta t)$ time steps, so the heat delivered
downstream is the heat injected $k_p$ steps earlier less the corresponding losses:
%
\begin{equation}
  Q_{p,t}^{\mathrm{delivered}} =
    Q_{p,\,t-k_p} - Q_{p,\,t-k_p}^{\mathrm{loss,pipe}}
  \quad \forall\, p,\; t \geq k_p{+}1.
  \label{eq:delayed_balance}
\end{equation}
%
Isolating delay on its own rung is what allows it to be separated from the linearisation
error rather than confounded with it. It also produces a clean negative
result, and the reason is visible in \eqref{eq:delay_time}: on this network every trunk
travel time is under an hour (the longest is about 35 minutes), so $k_p = 0$ for every
pipe and \eqref{eq:delayed_balance} collapses to the undelayed balance. $\Lsix$ and
$\Lthree$ are consequently the same programme, with identical variable, binary and
constraint counts and an identical optimum to the last digit. The delay effect at hourly
resolution is exactly zero rather than merely small, and this is a statement about the
temporal resolution: at sub-hourly resolution, or on a network with kilometres more trunk,
$k_p$ would not vanish.
